\SourceLeaf{17}
constatque, $T,U,V$ resp: proportionales esse cosinibus angulorum rectae, quae superficiei normalis est, cum rectis axibus coordinatarum $x,y,z$ parallelis, i. e. angulorum $QX,QY,QZ$. Hinc quidem colligitur, esse
\[
\begin{aligned}
\cos QX&=\frac{\pm T}{\BareRoot{(TT+UU+VV)}}\\
\cos QY&=\frac{\pm U}{\BareRoot{(TT+UU+VV)}}\\
\cos QZ&=\frac{\pm V}{\BareRoot{(TT+UU+VV)}}
\end{aligned}
\]
sed ambiguum manet, utrum signa superiora, an inferiora adoptare oporteat. Quod ut decidamus, capiamus in recta $PQ$ superficiei in $P$ normali \textit{extrorsumque} directa punctum $P'$ ipsi $P$ infinite proximum, sitque distantia $PP'=\mathrm{d}w$. Erunt itaque coordinatae puncti $P'$ resp.
\[
\begin{aligned}
x+\mathrm{d}w.\cos QX&=x+\mathrm{d}x\\
y+\mathrm{d}w.\cos QY&=y+\mathrm{d}y\\
z+\mathrm{d}w.\cos QZ&=z+\mathrm{d}z
\end{aligned}
\]
adeoque incrementum valoris functionis $W$ inde a puncto $P$ (ubi est $=0$) usque ad punctum $P'$
\[
\begin{aligned}
&=\mathrm{d}w.(T\cos QX+U\cos QY+V\cos QZ)\\
&=\pm\mathrm{d}w.\BareRoot{(TT+UU+VV)}
\end{aligned}
\]
Hinc patet, signa superiora valere, si recedendo a corporis soliditate functio $W$ nanciscitur valorem positivum, et proin negativum ingrediendo corporis soliditatem, contra signa inferiora valere in casu opposito. Revera quum superficies nostra tum corporis soliditatem a reliquo spatio vacuo separet, tum spatii partes eas, ubi $W$ positivum valorem obtinet, ab iis, ubi valor functionis $W$ fit negativus, generaliter loquendo \textit{vel} extra corpus valor functionis $W$ positivus erit, intra negativus, in quo casu signa superiora accipienda erunt, \textit{vel} functio $W$ negativa erit extra, positiva intra corpus, in quo casu signa inferiora valebunt.

Cosinus angulorum reliquorum, quibus in formulis nostris opus est, adhuc facilius evolvuntur. Habemus scilicet
\[
\begin{aligned}
a&=x+r\cos MX\\
b&=y+r\sin MY\\
c&=z+r\sin MZ
\end{aligned}
\]
